% =========================================================================
% OTIMIZAÇÃO COMBINATÓRIA EM GRAFOS BIPARTIDOS
% =========================================================================
\chapter{Otimização Combinatória em Grafos Bipartidos}\label{sec:bipartidos}
O Problema do Fluxo Máximo é frequentemente utilizado como base para resolver outros problemas de otimização combinatória via redução polinomial (veja~\cite{cook1998combinatorial}). O processo consiste em mapear a instância de um problema original para uma rede de fluxo equivalente, computar o fluxo máximo e traduzir o estado da rede residual de volta para a solução do problema original.

Um digrafo $D = (V, A)$ é classificado como bipartido se o seu conjunto de vértices $V$ puder ser particionado em dois subconjuntos disjuntos independentes, convencionalmente denotados por $L$ (esquerdo) e $R$ (direito), de modo que todo arco $a \in A$ conecte obrigatoriamente um vértice de $L$ a um vértice de $R$. Essa estrutura permite modelar problemas de alocação.


\section{Emparelhamento Máximo}\label{sec:emparelhamento}

O Problema do Emparelhamento Bipartido de Cardinalidade Máxima (\textit{Maximum Bipartite Matching}) consiste em encontrar um subconjunto de arestas $M \subseteq A$ de cardinalidade máxima tal que nenhum vértice seja incidente a mais de uma aresta em $M$.

O problema reduz-se ao cômputo de fluxo máximo em uma rede direcionada $D' = (V', A')$ construída a partir do grafo bipartido $G = (L \cup R, E)$:
\begin{enumerate}
    \item Adicionam-se uma superfonte $s$ e um supersumidouro $t$.
    \item Para cada $u \in L$, insere-se o arco $(s, u)$ com capacidade $c(s,u) = 1$.
    \item Para cada $v \in R$, insere-se o arco $(v, t)$ com capacidade $c(v,t) = 1$.
    \item Para cada aresta $\{u,v\} \in E$ com $u \in L$ e $v \in R$, insere-se o arco $(u, v)$ com capacidade $c(u,v) = 1$ (ou $\infty$).
\end{enumerate}

Pelo Teorema da Integralidade (Teorema~\ref{teo:integralidade}), como todas as capacidades incidentes a $s$ e $t$ são unitárias, existe fluxo máximo inteiro onde $f(u,v) \in \{0, 1\}$ para todo arco. As restrições de capacidade nos arcos incidentes a $s$ e $t$ asseguram que cada vértice de $L$ e $R$ participe de no máximo uma unidade de fluxo. Logo, o valor do fluxo máximo $|f^*|$ é igual à cardinalidade do emparelhamento máximo $|M|$, e os arcos com $f(u,v) = 1$ definem os pares emparelhados.

\begin{figure}[!ht]
\centering
\begin{tikzpicture}[thick]
\begin{scope}[xshift=0cm]
  \tikzsubcaption{(1.5,3.8)}{(a) Grafo bipartido original}
  \foreach \i/\y in {1/3, 2/1.5, 3/0} {
    \node[bip L] (l\i) at (0,\y) {$l_\i$};
    \node[bip R] (r\i) at (3,\y) {$r_\i$};
  }
  \draw[-] (l1) -- (r1); \draw[-] (l1) -- (r2);
  \draw[-] (l2) -- (r2); \draw[-] (l2) -- (r3);
  \draw[-] (l3) -- (r1); \draw[-] (l3) -- (r3);
\end{scope}
\node at (4.8,1.5) {\Huge $\Rightarrow$};

\begin{scope}[xshift=6.2cm]
  \tikzsubcaption{(2.5,3.8)}{(b) Rede de fluxo equivalente $D'$}
  \node[source node] (s) at (0,1.5) {$s$};
  \node[sink node] (t) at (5,1.5) {$t$};
  \foreach \i/\y in {1/3, 2/1.5, 3/0} {
    \node[bip L,font=\footnotesize] (fl\i) at (1.8,\y) {$l_\i$};
    \node[bip R,font=\footnotesize] (fr\i) at (3.2,\y) {$r_\i$};
    \draw[flow edge] (s) -- node[above,font=\footnotesize] {$1$} (fl\i);
    \draw[flow edge] (fr\i) -- node[above,font=\footnotesize] {$1$} (t);
  }
  \draw[flow edge,blue,line width=1.5pt] (fl1) -- (fr1); \draw[flow edge] (fl1) -- (fr2);
  \draw[flow edge,blue,line width=1.5pt] (fl2) -- (fr2); \draw[flow edge] (fl2) -- (fr3);
  \draw[flow edge] (fl3) -- (fr1); \draw[flow edge,blue,line width=1.5pt] (fl3) -- (fr3);
\end{scope}
\end{tikzpicture}

\caption{Redução do Emparelhamento Bipartido a Fluxo Máximo. (a) Grafo bipartido original com $L = \{l_1,l_2,l_3\}$ e $R = \{r_1,r_2,r_3\}$. (b) Rede de fluxo $D'$: arcos de $s$ para $L$ e de $R$ para $t$ com capacidade~1; arcos de $L$ para $R$ herdados do grafo original. Arcos azuis indicam o emparelhamento máximo obtido com $|f^*|=3$.}
\label{fig:fluxo_bipartido}
\end{figure}


\section{Cobertura de Vértices Mínima e o Teorema de König}\label{sec:cobertura}

O Problema da Cobertura de Vértices Mínima (\textit{Minimum Vertex Cover} - MVC) consiste em encontrar um subconjunto de vértices $C \subseteq V$ de cardinalidade mínima tal que toda aresta de $E$ seja incidente a pelo menos um vértice de $C$. Em grafos gerais o problema é $\mathcal{NP}$-difícil, mas em grafos bipartidos admite resolução em tempo polinomial.

A relação de dualidade é expressa pelo Teorema de König~\cite{korte2018combinatorial}:

\begin{teorema}[König, 1931][teo:konig]
Em qualquer grafo bipartido $G = (V, E)$, a cardinalidade de um emparelhamento máximo é igual à cardinalidade de uma cobertura de vértices mínima:
\begin{equation}
    \max |M| = \min |C|
\end{equation}
\end{teorema}

\intuicao{Dualidade Min-Max} Como cada aresta de um emparelhamento $M$ exige ao menos um vértice distinto para ser coberta, vale $|C| \ge |M|$ para qualquer grafo. Em grafos bipartidos, a igualdade é atingida pelo corte mínimo associado à rede de fluxo.

\begin{demonstracao}
O teorema decorre diretamente do Teorema do Fluxo Máximo e Corte Mínimo (Teorema~\ref{teo:mfmc}) aplicado à rede $D'$.

Após computar o fluxo máximo em $D'$, executa-se uma busca a partir de $s$ no digrafo residual $D'_f$ percorrendo arcos com $c_f > 0$. Seja $S$ o conjunto de vértices alcançáveis e $T = V' \setminus S$ os inalcançáveis, definindo um corte mínimo $(S, T)$.

Define-se $C = (L \cap T) \cup (R \cap S)$.

Como $(S, T)$ é um corte mínimo, não há arcos residuais com capacidade positiva de $S$ para $T$. Isso implica que não existem arestas de $L \cap S$ para $R \cap T$ no grafo original; logo, toda aresta tem extremidade em $L \cap T$ ou em $R \cap S$, de modo que $C$ é uma cobertura válida de vértices. Além disso, a capacidade do corte $(S, T)$ na rede é exatamente $|L \cap T| + |R \cap S| = |C|$. Pelo Teorema~\ref{teo:mfmc}, $\ccap(S, T) = |f^*| = |M^*|$, provando que $|C| = |M^*|$.
\end{demonstracao}

\implicacoes{e Complexidade} A redução ao fluxo máximo permite resolver o problema da cobertura de vértices mínima em grafos bipartidos em tempo $\mathcal{O}(|E|\sqrt{|V|})$ via Dinic ou Hopcroft-Karp.

\begin{figure}[!ht]
\centering
\begin{tikzpicture}[thick]
  \node[bip label L] at (2,4.5) {$\mathbf{L}$};
  \node[bip label R] at (4,4.5) {$\mathbf{R}$};

  \node[source node] (s) at (0,2) {$s$};
  \node[sink node] (t) at (6,2) {$t$};

  \node[source node] (l1) at (2,3.5) {$l_1$};
  \node[sink node,cover node] (l2) at (2,2) {$l_2$};
  \node[source node] (l3) at (2,0.5) {$l_3$};

  \node[sink node] (r1) at (4,3.5) {$r_1$};
  \node[sink node] (r2) at (4,2) {$r_2$};
  \node[source node,cover node] (r3) at (4,0.5) {$r_3$};

  \draw[->] (s) -- (l1);
  \draw[sat edge] (s) -- node[above,font=\footnotesize\bfseries,text=cutviolet!90!black,sloped]{corte} (l2);
  \draw[->] (s) -- (l3);

  \draw[match edge,->] (l1) -- (r3);
  \draw[match edge,->] (l2) -- (r1);
  \draw[->] (l2) -- (r2);
  \draw[->] (l3) -- (r3);

  \draw[->] (r1) -- (t);
  \draw[->] (r2) -- (t);
  \draw[sat edge] (r3) -- node[above,font=\footnotesize\bfseries,text=cutviolet!90!black,sloped]{corte} (t);

  \node[legend node,fill=green!20] (legS) at (0, -1.0) {};
  \node[right,font=\footnotesize] at (legS.east) {$\in S$ (Alcançável no residual)};

  \node[legend node,fill=red!20] (legT) at (0, -1.7) {};
  \node[right,font=\footnotesize] at (legT.east) {$\in T$ (Inalcançável)};

  \node[legend node,cover node] (legC) at (0, -2.4) {};
  \node[right,font=\footnotesize] at (legC.east) {$\in C$ (Cobertura Mínima)};

  \node[font=\footnotesize,align=center] at (3, -3.1) {\textbf{Resultado:} $C = \{l_2, r_3\}$, $|C| = |M| = 2$};

\end{tikzpicture}
\caption{Teorema de König: partição de vértices em $S$ e $T$ via corte mínimo residual. A cobertura mínima de vértices é $C = (L\cap T) \cup (R \cap S) = \{l_2, r_3\}$, com $|C| = |M| = 2$.}
\label{fig:konig_corte}
\end{figure}


\section{Conjunto Independente Máximo e o Teorema de Gallai}\label{sec:conjunto_independente}

O Problema do Conjunto Independente Máximo (\textit{Maximum Independent Set} - MIS) consiste em encontrar um subconjunto $I \subseteq V$ de cardinalidade máxima tal que nenhuma aresta una dois vértices de $I$.

Como um subconjunto $I \subseteq V$ é independente se, e somente se, o complementar $V \setminus I$ é uma cobertura de vértices, tem-se a relação direta:
\begin{equation}
    I = V \setminus C
\end{equation}
onde $C$ é a cobertura de vértices mínima obtida pelo Teorema de König. Portanto:
\begin{equation}
    |I^*| = |V| - |C^*| = |V| - |M^*|
\end{equation}
A extração de $I$ opera em tempo $\mathcal{O}(|V|)$ a partir de $C$.

\begin{figure}[!ht]
\centering
\begin{tikzpicture}[thick]
  \begin{scope}[shift={(0,0)}]
    \tikzsubcaption{(2.5,4.2)}{(a) Cobertura de Vértices $C$}

    \node[bip L] (l1) at (1.7,3.5) {$l_1$};
    \node[bip L] (l3) at (1.7,0.5) {$l_3$};
    \node[bip R] (r1) at (3.3,3.5) {$r_1$};
    \node[bip R] (r2) at (3.3,2) {$r_2$};

    \node[bip L,fill=blue!40,cover node] (l2) at (1.7,2) {$l_2$};
    \node[bip R,fill=orange!40,cover node] (r3) at (3.3,0.5) {$r_3$};

    \draw[-] (l1) -- (r3);
    \draw[-] (l2) -- (r2);
    \draw[-] (l3) -- (r3);
    \draw[-] (l2) -- (r1);

    \node[font=\footnotesize,below] at (2.5,-0.2) {$C = \{l_2, r_3\}$};
  \end{scope}

  \begin{scope}[shift={(6,0)}]
    \tikzsubcaption{(2.5,4.2)}{(b) Conjunto Independente $I = V \setminus C$}

    \node[bip L,fill=blue!40,cover node] (l1) at (1.7,3.5) {$l_1$};
    \node[bip L,fill=blue!40,cover node] (l3) at (1.7,0.5) {$l_3$};
    \node[bip R,fill=orange!40,cover node] (r1) at (3.3,3.5) {$r_1$};
    \node[bip R,fill=orange!40,cover node] (r2) at (3.3,2) {$r_2$};

    \node[bip L] (l2) at (1.7,2) {$l_2$};
    \node[bip R] (r3) at (3.3,0.5) {$r_3$};

    \draw[-] (l1) -- (r3);
    \draw[-] (l2) -- (r2);
    \draw[-] (l3) -- (r3);
    \draw[-] (l2) -- (r1);

    \node[font=\footnotesize,below] at (2.5,-0.2) {$I = \{l_1, l_3, r_1, r_2\}$};
  \end{scope}
\end{tikzpicture}
\caption{Complementaridade entre Cobertura de Vértices e Conjunto Independente Máximo: $I = V \setminus C$.}
\label{fig:conjunto_independente}
\end{figure}


\section{Cobertura de Arestas Mínima}\label{sec:cobertura_arcos}

A \textbf{Cobertura de Arestas Mínima} (\textit{Minimum Edge Cover} - MEC) consiste no menor subconjunto $F \subseteq E$ tal que todo vértice em $V$ é incidente a pelo menos uma aresta de $F$.

A relação com o emparelhamento máximo é expressa pelas Identidades de Gallai~\cite{bondy2008graph}:

\begin{teorema}[Identidades de Gallai, 1959][teo:gallai]
Em qualquer grafo simples $G = (V, E)$ sem vértices isolados:
\begin{align}
    |M^*| + |C_V^*| &= |V| \label{eq:gallai_v} \\
    |M^*| + |F^*| &= |V| \label{eq:gallai_e}
\end{align}
onde $|M^*|$ é o tamanho do emparelhamento máximo, $|C_V^*|$ da cobertura de vértices mínima e $|F^*|$ da cobertura de arestas mínima.
\end{teorema}

\intuicao{Combinatória} Cada aresta de um emparelhamento cobre exatamente dois vértices distintos. Os $|V| - 2|M^*|$ vértices desemparelhados requerem cada um ao menos uma aresta adicional para cobertura, resultando em $|F^*| = |M^*| + (|V| - 2|M^*|) = |V| - |M^*|$.

\begin{demonstracao}
Dado um emparelhamento máximo $M^*$, para cada um dos $|V| - 2|M^*|$ vértices não cobertos seleciona-se uma aresta incidente (possível pois não há vértices isolados). O conjunto resultante $F$ cobre todos os vértices e satisfaz $|F| \le |M^*| + |V| - 2|M^*| = |V| - |M^*|$, donde $|F^*| \le |V| - |M^*|$. Reciprocamente, qualquer cobertura mínima $F^*$ gera uma floresta de estrelas com cardinalidade $|F^*| \ge |V| - |M^*|$. Logo, $|M^*| + |F^*| = |V|$. A equação~\eqref{eq:gallai_v} decorre de $|C_V^*| = |M^*|$ por König em grafos bipartidos.
\end{demonstracao}

A cobertura mínima de arestas $F^*$ é obtida a partir de $M^*$ em tempo $\mathcal{O}(|V| + |E|)$, completando o conjunto de reduções.


\section{O Quadro Unificado das Quatro Dualidades (König e Gallai)}

As relações entre as quatro grandezas combinatórias em grafos bipartidos estão resumidas na Tabela~\ref{tab:dualidades}:

\begin{table}[!ht]
\centering
\small
\setlength{\tabcolsep}{12pt}
\begin{tabular}{l l l}
\toprule
\multicolumn{1}{c}{\textbf{Objetivo}} & \multicolumn{1}{c}{\textbf{Domínio de Vértices ($V$)}} & \multicolumn{1}{c}{\textbf{Domínio de Arestas ($E$)}} \\
\midrule
\textbf{Maximizar} & \hyperref[sec:conjunto_independente]{Conjunto Independente} $\alpha(G) = |I^*|$ & \hyperref[sec:emparelhamento]{Emparelhamento Máximo} $\nu(G) = |M^*|$ \\
\addlinespace[0.4em]
\textbf{Minimizar} & \hyperref[sec:cobertura]{Cobertura de Vértices} $\tau(G) = |C_V^*|$ & \hyperref[sec:cobertura_arcos]{Cobertura de Arestas} $\rho(G) = |F^*|$ \\
\bottomrule
\end{tabular}
\caption{Quatro dualidades combinatórias em grafos bipartidos: $\nu(G) = \tau(G)$ (König) e $\alpha(G) + \tau(G) = \nu(G) + \rho(G) = |V|$ (Gallai).}
\label{tab:dualidades}
\end{table}

Combinando a equivalência de König $\nu(G) = \tau(G)$ com Gallai, decorre a identidade:
\begin{equation}
    \alpha(G) = |V| - \tau(G) = |V| - \nu(G) = \rho(G)
\end{equation}
mostrando que a cardinalidade do conjunto independente máximo coincide com a da cobertura de arestas mínima em grafos bipartidos sem vértices isolados.

\begin{figure}[!ht]
\centering
\begin{tikzpicture}[thick]
  \node[concept box, fill=blue!10, text width=3.5cm] (mis) at (0, 6) {\textbf{Conjunto Independente} \\ MIS \\ (Maximizar, Vértices)};

  \node[concept box, fill=orange!10, text width=3.5cm] (match) at (9, 6) {\textbf{Emparelhamento} \\ Matching \\ (Maximizar, Arestas)};

  \node[concept box, fill=blue!10, text width=3.5cm] (mvc) at (0, 0) {\textbf{Cobertura de Vértices} \\ MVC \\ (Minimizar, Vértices)};

  \node[concept box, fill=orange!10, text width=3.5cm] (mec) at (9, 0) {\textbf{Cobertura de Arestas} \\ MEC \\ (Minimizar, Arestas)};

  \draw[<->] (mis) -- node[left, font=\footnotesize, align=right] {Complemento\\$|I| = |V| - |C_V|$} (mvc);
  \draw[<->] (match) -- node[right, font=\footnotesize, align=left] {Gallai\\$|M| + |F| = |V|$} (mec);

  \draw[<->] (match) -- node[sloped, above, font=\footnotesize, pos=0.25] {König} node[sloped, below, font=\footnotesize, pos=0.25] {$|M| = |C_V|$} (mvc);
  \draw[<->, dashed] (mis) -- node[sloped, above, font=\footnotesize, pos=0.25] {Gallai (Derivado)} node[sloped, below, font=\footnotesize, pos=0.25] {$|I| = |F|$} (mec);
\end{tikzpicture}
\caption{Quatro grandezas duais em grafos bipartidos conectadas pelos teoremas de König e Gallai.}
\label{fig:quatro_dualidades}
\end{figure}


\section{O Teorema do Casamento de Hall}\label{sec:hall}

Seja $G = (L \cup R, E)$ um grafo bipartido. Um emparelhamento $M$ é $L$-saturante se cobre todos os vértices de $L$ ($|M| = |L|$). Para $S \subseteq L$, denota-se por $N(S) = \{v \in R \mid \exists\, u \in S, \, \{u,v\} \in E\}$ a vizinhança de $S$ em $R$.

\begin{teorema}[Casamento de Hall, 1935][teo:hall]
Existe um emparelhamento $L$-saturante em $G = (L \cup R, E)$ se, e somente se:
\begin{equation}
    |N(S)| \ge |S|, \quad \forall\, S \subseteq L
\end{equation}
\end{teorema}

\begin{demonstracao}
$(\Rightarrow)$ Se existe emparelhamento $L$-saturante $M$, as arestas de $M$ incidentes a $S$ associam cada vértice de $S$ a um vizinho distinto em $N(S)$. Logo, $|N(S)| \ge |S|$.

\medskip
\noindent $(\Leftarrow)$ Suponha por contradição que $|N(S)| \ge |S|$ para todo $S \subseteq L$, mas o emparelhamento máximo satisfaça $|M^*| < |L|$.

Pelo Teorema de König (Teorema~\ref{teo:konig}), existe cobertura de vértices $C$ com $|C| = |M^*| < |L|$. Sejam $C_L = C \cap L$ e $C_R = C \cap R$, com $|C| = |C_L| + |C_R|$.

Defina $S = L \setminus C_L$. Como $C$ cobre todas as arestas, não há arestas entre $S$ e $R \setminus C_R$, implicando $N(S) \subseteq C_R$ e, portanto, $|N(S)| \le |C_R|$.

Por hipótese de Hall, $|N(S)| \ge |S| = |L| - |C_L|$. Segue que:
$$|C_R| \ge |N(S)| \ge |L| - |C_L| \implies |C| = |C_L| + |C_R| \ge |L|$$
o que contradiz $|C| < |L|$. Logo, $|M^*| = |L|$.
\end{demonstracao}

A verificação computacional da condição de Hall e a construção do emparelhamento correspondente realizam-se em tempo $\mathcal{O}(|E|\sqrt{|V|})$ via fluxo máximo.

\begin{figure}[!ht]
\centering
\begin{tikzpicture}[thick]
    \node[bip L] (l1) at (0,3) {$l_1$};
    \node[bip L] (l2) at (0,1.5) {$l_2$};
    \node[bip L] (l3) at (0,0) {$l_3$};

    \node[bip R] (r1) at (3,3) {$r_1$};
    \node[bip R] (r2) at (3,1.5) {$r_2$};
    \node[bip R] (r3) at (3,0) {$r_3$};

    \draw[match edge] (l1) -- (r1);
    \draw[match edge] (l2) -- (r2);
    \draw[match edge] (l3) -- (r3);

    \draw[unmatch edge] (l1) -- (r2);
    \draw[unmatch edge] (l2) -- (r3);
    \draw[unmatch edge] (l3) -- (r1);

    \node[bip label L] at (0,3.7) {$\mathbf{L}$};
    \node[bip label R] at (3,3.7) {$\mathbf{R}$};
\end{tikzpicture}
\caption{Emparelhamento $L$-saturante (arestas azuis) satisfazendo a condição de Hall $|N(S)| \ge |S|$ para todo $S \subseteq L$.}
\label{fig:hall_casamento}
\end{figure}


\section{Equivalência com Algoritmos Dedicados e Simplificação Residual}

Algoritmos dedicados para emparelhamento bipartido, como Kuhn~\cite{kuhn1955hungarian} e Hopcroft-Karp~\cite{korte2018combinatorial}, correspondem a especializações dos métodos de Ford-Fulkerson e Dinic sob capacidades unitárias:

\begin{enumerate}
    \item \textbf{Capacidades Unitárias ($c \in \{0, 1\}$):} Cada aresta assume apenas dois estados: livre ($f = 0$, sentido $L \to R$) ou emparelhada ($f = 1$, sentido residual $R \to L$).
    \item \textbf{Arcos Reversos Determinísticos:} Como cada vértice $v \in R$ emparelha-se com no máximo um $u \in L$, o arco reverso no digrafo residual é único e dado por \lstinline{match[v]}.
    \item \textbf{Caminhos Alternantes:} Pela bipartição, todo caminho residual aumentante alterna arestas não emparelhadas e emparelhadas.
\end{enumerate}

Essa correspondência traduz-se em paridade assintótica:
\begin{itemize}
    \item \textbf{Kuhn vs. Ford-Fulkerson:} Executa DFS para encontrar um caminho alternante por iteração em tempo $\mathcal{O}(|V| \cdot |E|)$.
    \item \textbf{Hopcroft-Karp vs. Dinic:} Opera em fases com BFS global seguida de DFS bloqueadora. O número de fases é limitado por $\mathcal{O}(\sqrt{|V|})$, resultando em $\mathcal{O}(|E| \sqrt{|V|})$~\cite{korte2018combinatorial}, equivalente a Dinic sobre a rede unitária.
\end{itemize}

Para o caso ponderado (atribuição de custo mínimo), o Algoritmo Húngaro~\cite{kuhn1955hungarian} equivale ao método \textit{Successive Shortest Path} com potenciais nodais, analisado na Seção~\ref{sec:ssp}.


